% ============================================================================
% Section 1: Introduction (standalone)
% ============================================================================

\section{Introduction}\label{sec:introduction}

The study of Artificial Intelligence (AI) systems is undergoing a fundamental shift. Where language models were once evaluated on their ability to quickly produce fluent text in response to human prompts, the frontier has moved toward sustained, autonomous interaction within dynamic virtual and physical environments. AI systems are now expected to write and debug code across entire repositories, operate robotic platforms, and navigate virtual worlds, tasks that move beyond the demands of pretrained knowledge and extend towards the ability to plan, adapt, and progress over many hours of continuous execution.

Partial observability, long-horizon planning, and multi-modal perception are fundamental problems at the core of these emerging workflows. In these settings, agents are not only required to iterate on decisions that optimize for the successful completion of near-horizon tasks, but must also effectively plan for and maintain coherence across many thousands of environment interactions in pursuit of larger, often more abstract goals. In light of this new paradigm, there have been recent efforts to test the capabilities of frontier Large Language Models (LLMs) and the agent systems they underpin. Yet few existing benchmarks stress all three dimensions simultaneously under realistic conditions. Standard testbeds tend to isolate one axis: imperfect-information games emphasize equilibrium computation in short episodes, while open-ended environments test exploration but lack adversarial opponents or extended temporal commitments. As frontier language and vision-language models (VLMs) are increasingly deployed as autonomous agents, the gap between what existing benchmarks measure and what real-world autonomy demands has grown acute.

Pok\'{e}mon role-playing games (RPGs) offer a compelling case-study for this exact purpose. They combine a vast, partially observable state space with long-horizon, often non-linear narrative progression that demands thousands of cumulative decisions spanning exploration, resource management, combat, and spatial navigation. The recent surge of interest in Pok\'{e}mon-playing AI systems (Claude Plays Pok\'{e}mon~\citep{anthropic2025visible}, Gemini Plays Pok\'{e}mon~\citep{zhang2025geminiplayspokemon}, and GPT-5's completion of Pok\'{e}mon Red~\citep{engadget2025gptplayspokemon}) has demonstrated both the appeal of this domain and its limitations as a research vehicle. Fragmented across different games, harnesses, and evaluation criteria, the disconnected nature of these projects has made it difficult to rigorously disentangle both quantitative and qualitative metrics for success from the performance of the underlying model, sophistication of the harness architecture, or hardcoded assumptions about access to environment state.

This work addresses the fragmentation problem, advancing the state-of-the-art (SOTA) across three complementary dimensions. We first introduce an open \textbf{speedrunning benchmark}, as a robust foundation for evaluating frontier model and harness capabilities in long-duration, embodied settings and builds upon the testing framework established by the Speedrunning Track of the \pokeagent{} Challenge at NeurIPS 2025~\citep{karten2025pokeagent}. Our benchmark formalizes Pok\'{e}mon Emerald gameplay as a structured evaluation with standardized milestones, quantitative metrics, and infrastructure for fair cross-paradigm comparison. We have since extended it from the tutorial segment (through the first gym) to a long-horizon evaluation spanning many more segments of the game, where experiments proceed on the order of days and tens of thousands of environment interactions. Next, we present \textit{\pokeagent{},} a SOTA \textbf{custom VLM-agent harness} tasked with maintaining long-context coherence, developing and iterating upon plans, and effective problem solving across the Pok\'{e}mon Emerald storyline. Here, we detail the many insights spanning the development of \pokeagent{} from a simple four-module scaffold during the NeurIPS competition into a multi-agent system with abstractions for processing global session-level context, hierarchical memory management, skill-learning, subagent orchestration, and sandboxed code execution. We further compare PokeAgent against common CLI-agent harnesses (Claude Code, Codex CLI, Gemini CLI) revealing that, while coding-agent architectures are impressive out-of-the-box, they
nonetheless struggle to properly adapt to the unique set of constraints imposed by our speedrunning benchmark. Finally, we introduce \textbf{AutoEvolve}, a meta-harness framework and reset-free evolutionary algorithm designed to automate experiential learning by iteratively building agent scaffolding through in-context optimization over global trajectory segment history. Beginning from the simplest initialization, AutoEvolve seeks to bootstrap the development of its own harness by pruning and refining memories, distilling agent progress across entire episodes of environment interactions into skill and subagent abstractions, and refining system prompts accordingly with respect to the evolving harness. Unlike prior methods for model-level and agent self-adaptation that require full episodes and environment resets~\citep{agrawal2025gepa}, AutoEvolve operates in a fully online, reset-free manner, targeting continuous self-improvement in settings where environment resets are expensive or unfeasible. 

The remainder of this work is organized as follows. Chapter~\ref{sec:background} provides the necessary background, introducing relevant foundations in Markov Decision Processes, the Transformer architecture, reinforcement learning, and the Model Context Protocol, before surveying the broader motivation for this work. Chapter~\ref{sec:related_work} surveys related work across experiential learning, long-horizon benchmarks, embodied agents in RPG settings, and self-adaptive methods. Chapter~\ref{sec:formalization} formalizes our conception of \textbf{Embodied Agent Environments (EAEs)} and \textbf{Agentic Harnesses}, including a spectrum of harness configurations that frames the experimental comparisons in later chapters. Chapter~\ref{sec:benchmark} describes the speedrunning benchmark in detail: its environment design, evaluation criteria, and baselines. Chapter~\ref{sec:pokeagent} traces the development of the \pokeagent{} harness from its competition-era origins through its current multi-agent architecture, presenting early performance results that motivate our interest in automated harness evolution. Chapter~\ref{sec:autoevolve} presents AutoEvolve, including its methodology, experimental design across many agent-evaluation settings, and a discussion of our results. Chapter~\ref{sec:outlook} concludes with an exploration of the implications of this work and fruitful areas for future investigation across model distillation, ultra-long context (ULC) evaluations, and extensions to domains beyond Pok\'{e}mon.
